%======================================================================
\section{Setting and standing hypotheses}\label{sec:setting}
%======================================================================
$Y$ is an integer homology sphere, $K\subset Y$ a knot, $X=Y\setminus\nu K$, with meridian $\mu$ and
Seifert longitude $\lambda$ on $\partial X=T^2$.  $R(X)=\Hom(\pi_1X,\SU(2))/\mathrm{conj}$, and $R^{\rm irr}(X)$ is
its irreducible part (an open subset).  For $\rho\in R(X)$ choose a conjugate with
$\rho(\mu)=\mathrm{diag}(e^{2\pi ia},e^{-2\pi ia})$, $\rho(\lambda)=\mathrm{diag}(e^{2\pi ib},e^{-2\pi ib})$; the pair
$(a,b)\in T^2=\R^2/\Z^2$ is defined up to $(a,b)\mapsto(-a,-b)$, and
\[
r\colon R(X)\to P:=T^2/\pm,\qquad \rho\mapsto[(a,b)],
\]
is the restriction to the boundary.  $P$ is the pillowcase, a sphere with four orbifold points (the corners),
the images of the central representations $\rho(\mu),\rho(\lambda)\in\{\pm1\}$.  We write $\Gam:=r(R^{\rm irr}(X))$.

Irreducible instanton homology: for an integer homology sphere $Z$, $\CI_*(Z)$ denotes Floer's complex generated by
the irreducible flat connections (after a small holonomy perturbation), $\Z/8$-graded, with its Chern--Simons
filtration.  We work over $\F=\Z/2$ unless stated otherwise; orientations are a separate issue (\S\ref{sec:workflow}).

\subsection{Hypotheses}
\begin{enumerate}[label=(H\arabic*),leftmargin=2.6em]
\item\label{H:alex} $\Delta_K(e^{2i\theta})\neq0$ for all $\theta$.
\item\label{H:smooth} $R^{\rm irr}(X)$ is a smooth $1$-manifold, and $H^1(X;\mathrm{ad}\,\rho)$ is $1$-dimensional for every
$\rho\in R^{\rm irr}(X)$ (the Zariski tangent space equals the tangent space).
\end{enumerate}
For $Y\ne S^3$ we add
\begin{enumerate}[label=(H\arabic*),leftmargin=2.6em,resume]
\item\label{H:Y} every irreducible flat connection $\rho$ on $Y$ is nondegenerate and has $\rho(K)\ne\pm1$.
\end{enumerate}

\subsection{Consequences}
\begin{proposition}\label{prop:basic}
Assume \ref{H:alex}--\ref{H:smooth}.
\begin{enumerate}[label=(\alph*)]
\item $R^{\rm irr}(X)$ is compact: a finite disjoint union of circles $\gamma_1,\dots,\gamma_m$.  Its image $\Gam$ is a
finite union of immersed closed curves in $P$: $r$ is an immersion.
\item If $Y=S^3$, $\Gam$ avoids the two edges $a\in\{0,\tfrac12\}$ of $P$, in particular the corners.  In general, $\Gam$ meets the corners only at representations
with $\rho(\mu)=\pm1$; those with $\rho(\mu)=1$ are the irreducible flat connections on $Y$.
\item $R(X)$ is invariant under $\chi\colon\rho\mapsto\rho\cdot\chi$, where $\chi\colon\pi_1X\to\{\pm1\}$ is the
nontrivial character; on $T^2$ this is $(a,b)\mapsto(a+\tfrac12,b)$.  Hence $\Gam$ is invariant under the
translation $a\mapsto a+\tfrac12$.
\end{enumerate}
\end{proposition}
\begin{proof}
(a) By Klassen's theorem (and its extension to homology spheres), an abelian representation with
$\rho(\mu)=e^{i\theta}$ is a limit of irreducibles only if $\Delta_K(e^{2i\theta})=0$.  Under \ref{H:alex} no abelian
representation is such a limit, so $R^{\rm irr}(X)$ is closed in the compact space $R(X)$.  With \ref{H:smooth} it is a
compact $1$-manifold.  The image of $dr\colon H^1(X;\mathrm{ad}\,\rho)\to H^1(T^2;\mathrm{ad}\,\rho)$ is Lagrangian for
the cup-product symplectic form, hence $1$-dimensional; with $\dim H^1(X;\mathrm{ad}\,\rho)=1$, $dr$ is injective.
(b) If $Y=S^3$ then $\mu$ normally generates $\pi_1X$, so $\rho(\mu)=\pm1$ forces $\rho$ central.  In general
$\rho(\mu)=1$ means $\rho$ factors through $\pi_1Y$.  (c) $\lambda$ lies in the commutator subgroup, so $\chi(\lambda)=1$.
\end{proof}

\subsection{Surgery lines}
For a slope $p/q$ (with $p,q$ coprime, $p/q=1/0$ for $Y$ itself) put
\[
L_{p/q}:=\{pa+qb\equiv0\},\qquad L^w_{p/q}:=\{pa+qb\equiv\tfrac12\}\subset T^2,
\]
and their images in $P$.  Flat $\SU(2)$ connections on $Y_{p/q}(K)$ restrict to representations of $\pi_1X$ with
$\rho(\mu^p\lambda^q)=1$, i.e.\ to points of $L_{p/q}$.  For $p$ even, $L^w_{p/q}$ describes the flat connections on
the nontrivial $\SO(3)$ bundle; in particular $L^w_{0}=\{b=\tfrac12\}$ describes $I(Y_0(K),w)$ with $w$ the
nontrivial bundle (admissible since $Y_0$ has $b_1=1$).

\begin{proposition}[generators]\label{prop:gens}
Assume \ref{H:alex}--\ref{H:smooth}, and $p\neq0$.
\begin{enumerate}[label=(\alph*)]
\item The irreducible flat connections on $Y_{p/q}(K)$ are in bijection with $r^{-1}(L_{p/q})\subset R^{\rm irr}(X)$.
The only reducible flat connection on $Y_{1/n}(K)$ is the trivial one.
\item Such a flat connection $\rho$ is nondegenerate ($H^1(Y_{p/q};\mathrm{ad}\,\rho)=0$) if and only if $\Gam$ is
transverse to $L_{p/q}$ at $r(\rho)$.
\item The flat connections on $(Y_0(K),w)$ are in bijection with $r^{-1}(L^w_0)$, all irreducible; nondegeneracy is
as in (b).
\end{enumerate}
\end{proposition}
\begin{proof}
(a) A representation of $\pi_1X$ with $\rho(\mu^p\lambda^q)=1$ extends uniquely over the solid torus, and
$\pi_1X\to\pi_1Y_{p/q}$ is onto, so irreducibility is preserved both ways.  An abelian $\rho$ has $\rho(\lambda)=1$, so
$\rho(\mu)^p=1$; for $p=\pm1$ this is trivial.  (b) Mayer--Vietoris for $\mathrm{ad}\,\rho$ along $T^2$: since
$\rho|_{T^2}$ is non-central abelian (by Proposition~\ref{prop:basic}(b) and \ref{H:Y}),
$H^0(T^2)=\R=H^0(\text{solid torus})$ and $H^1(T^2;\mathrm{ad}\,\rho)=T_{(a,b)}T^2$ is $2$-dimensional.  Then
$H^1(Y_{p/q};\mathrm{ad}\,\rho)=\ker\big(H^1(X)\oplus H^1(S^1\times D^2)\to H^1(T^2)\big)$, with $H^1(X)=T_\rho R^{\rm irr}(X)$
mapping injectively by $dr$ and $H^1(S^1\times D^2)=\R$ mapping onto the tangent line of $L_{p/q}$.  The kernel
vanishes iff the two lines are transverse.  (c) is the same with $\rho(\lambda)=-1$; abelian representations have
$\rho(\lambda)=1$, so none occur.
\end{proof}

\begin{remark}
Transversality can be achieved for every slope by an arbitrarily small holonomy perturbation supported in a collar
of $T^2$; in the pillowcase this replaces $L_{p/q}$ by a $C^1$-close curve $L'_{p/q}$ (Herald's collar
perturbations).  The regularity of instanton moduli spaces may require further small perturbations.  All statements
below are stable under such perturbations: the generator data move by $O(\epsilon)$ and the filtered homotopy type
is taken in the limit, as in the definition of Daemi's $\Gamma_Y$.
\end{remark}
